\documentclass[11pt]{article}
\usepackage{amsmath,amssymb,amsthm,hyperref}
\newtheorem{theorem}{Theorem}
\newtheorem{proposition}[theorem]{Proposition}
\newcommand{\E}{\mathcal E}
\title{Variable-arity feedback through scope-atomic causal frontier strategies}
\date{16 November 2026}
\begin{document}\maketitle
\paragraph{Audit and status.}
The written construction and proofs have passed independent review at the stated scope. No proof-assistant certification or TeX compilation is claimed. This is a parameterized construction, not an unrestricted shared polynomial contentad or a completion of O1. The input interface declares that one entire active frontier is an atomic sampling scope. Conditional joint laws within that scope are arbitrary couplings of the prescribed primitive rows; different rounds use causal conditional kernels. Cuts retain the complete raw past and cannot split a scope. These are primitive, checkable rules, not admission defined as preservation under all desired composites. The construction accommodates the full variable-arity terminating-tree grade $R_U C$, proves a stopping-frontier cut/assembly theorem, and retains native coefficients and unused grades. Its failure boundary is a proved obstruction to arbitrary scope-splitting cuts. Global mutable state and arbitrary shared effects remain outside the theorem.

\section{Independent inputs and native categories}
Fix a field $k$. The original O1 augmented context rig
\[
 (\mathbb R_{\geq0})_{\mathrm{disc}}\times\mathrm{FinVect}_k
\]
uses addition/direct sum and multiplication/tensor. A finite normalized row is a finite raw set $W$, weights $p_w\geq0$ summing to one, and augmented vector spaces $E_w\to k$. A native arrow is an actual raw function $r$ with $p'_{w'}=\sum_{r(w)=w'}p_w$, together with actual augmentation-preserving arrows $E_w\to E'_{r(w)}$. This is the exact coproduct decomposition of the original aggregation arrows, not an identification of coefficients with their carriers. Write $\E=\mathrm{FinVect}_k/k$, with its categorical finite products and terminal object $k\xrightarrow{1}k$. The same arguments work for an augmented coefficient category with the indicated finite products and augmentation-induced tensor-to-product maps.

For the terminating output only, use countable raw sets, weights of total at most one, and the same individual augmented coefficients. Native arrows still require exact fibre pushforward and actual coefficient arrows. This is a precisely declared subprobability enlargement, not a normalized output claim. It can be realized in the extended scalar/vector rig $([0,\infty],\leq)\times\mathrm{Vect}_k$, but the arrows considered here are the exact-pushforward subcategory. Nothing below identifies absent mass with a specific divergent trajectory.

\paragraph{The product-coefficient dependent bind is constructed directly.}
For a countable subprobability context $(W,p,E)$ and contexts $(Z_w,q_w,F_w)$, define the bound raw set as $\coprod_w Z_w$, weights $p_wq_w(z)$, and coefficient $E_w\times F_{w,z}$. On maps, a base arrow $r$ and compatible fibre arrows $s_w:Z_w\to Z'_{r(w)}$ give $(w,z)\mapsto(r(w),s_w(z))$ and the actual product of the displayed coefficient arrows. Exact pushforward follows by summing over each fibre and using first the fibre equations and then the base equation. These definitions require no distributivity of categorical products over vector-space direct sums. If an aggregation arrow is written in the scalar/vector rig, its vector part is the unique arrow out of the direct sum with the specified individual coefficient arrows. Identity and composition hold on these summands. The unit is a singleton of weight one with terminal coefficient. Associativity is the raw dependent-tuple bijection and canonical finite-product associator; its pentagon follows from projection uniqueness. This supplies precisely the dependent assembly operation used below without citing a rig theorem for a nonexistent distributive product structure on $\E$.

Independently fix a symmetric monoidal category $\mathcal D$ of content coefficients. A full ordered polynomial collection $C$ consists of all shapes $c\in C_n$ for every finite $n$, with coefficients $D_c\in\mathcal D$. Maps preserve the ordered positions and carry actual arrows $D_c\to D'_{\alpha(c)}$. No restriction to homogeneous arity is made. A primitive normalized process
\[
 f:A+U\longrightarrow F_C(B+U)
\]
has finite raw row $W_x$, law $p_x$, context coefficients $E_{x,w}$, and body assigning a shape $c(x,w)$ of arity $n$ and a list of $n$ outputs in $B+U$. Finite rows can be replaced by countable normalized rows throughout. Inputs and loop states may be arbitrary small sets; per external input the reachable finite raw histories remain countable. All raw zero witnesses are retained.

The global grade $R_U C$ is the entire collection of nonempty finite planar $C$-trees, with internal edges labelled in $U$, external ordered leaves, and nodewise tensor product of all content coefficients. It includes shapes that no process history realizes. The construction below never replaces this collection by the support or image of a body map.

\section{The independently specified scope/strategy interface}
A finite raw history begins at one unexecuted external root $a$. In one round, every currently active occurrence executes once. A witness at a node supplies its ordered children: $B$-children become permanently settled leaves and $U$-children become the next active frontier. Histories record all primitive raw witnesses, all addresses, and all inputs, including witnesses of weight zero. The frontier $I(h)$ is a finite ordered set of addresses. The primitive input $x_i(h)$ at each address is determined by the history, not by the current draw. A history is terminal exactly when $I(h)$ is empty. The root counts as the first round; nullary witnesses terminate without inventing a leaf.

A causal frontier strategy $\sigma$ supplies, at \emph{every syntactically possible nonterminal history}, a normalized law
\[
 q_h\text{ on }\prod_{i\in I(h)}W_{x_i(h)},\qquad
 (\pi_i)_*q_h=p_{x_i(h)}.                                      \tag{S}
\]
The condition is imposed even at zero-probability histories. Thus a strategy is a complete contingent table rather than its almost-sure equivalence class. No measurability issue is hidden: each reachable history set is countable and all current rows are discrete. Products are finite. At the initial singleton frontier, (S) fixes the root law. Subsequent frontiers can contain any finite number of occurrences with different inputs and different possible arities.

The distinction between randomness and coefficients is retained. A frontier tuple $\mathbf w$ has joint context coefficient
\[
 J(h,\mathbf w)=\prod_{i\in I(h)}E_{x_i(h),w_i}
\]
in $\E$, not the independent tensor. Its entire prior history remains available to select later $q_h$. A frontier may share a fair bit across arbitrarily many active occurrences; (S) does not require within-frontier independence. On the other hand it requires each primitive marginal to remain the declared row conditional on the \emph{whole prior history}. Persistent random state changing these conditional marginals must be explicitly included in the primitive input with its conditional rows. It is not supplied for free by the word ``sharing.''

This interface is local: testing (S) requires only marginalization of one finite or countable table. It neither enumerates future routings nor presupposes an assembled native comparison.

\paragraph{Consumption of the controlled original crossing.}
At a history $h$, the addressed frontier is already known before the draw. Apply the companion controlled crossing to this ordered occurrence list and the primitive context rows $W_{x_i(h)}$. Its control fibre is exactly the set of $q_h$ satisfying (S), and its actual coefficient is exactly $J(h,\mathbf w)$. Advancing the history applies the supplied primitive shape/output body to each coordinate, so new positions are determined only after that draw. Thus the evaluator consumes a succession of those local controls; it does not flatten arbitrary couplings of bound rows and then presume an inverse temporal interchange. The companion's nonanticipation equation is precisely the two-round marginal condition obtained by multiplying a prefix law by these conditional controls. Zero-prefix controls are retained here as part of $\sigma$, not reconstructed from the flattened law.

There is a coefficient-doctrine distinction. A source staged execution may use a tensor of the products for its successive rounds, whereas this target uses the product of all node coefficients. The directed comparison from the former to the latter is obtained, for each node, by augmenting all other round factors and then projecting to that node within its round. Product universality gives the map. All regroupings have those same node projections, proving compatibility with staged evaluation. It is generally noninvertible. Neither this bridge nor the old O1 injection identifies tensor and product coefficients or assumes a strong temporal interchange.

\section{Terminating recursive feedback at the full tree grade}
For a finite terminal history $t$, let $h_0,\ldots,h_{d-1}$ be its successive nonterminal prefixes and $\mathbf w_j$ its frontier draws. Set
\[
 P_\sigma(t)=\prod_{j<d}q_{h_j}(\mathbf w_j),\qquad
 E_t=\prod_{v\in\mathrm{nodes}(t)}E_{x_v,w_v}.                 \tag{T}
\]
The raw terminal history set contains all syntactically possible terminating histories, not merely positive ones. Breadth-first expansion gives a bijection with all compatible finite raw execution trees of the primitive process. Forgetting raw labels gives a body at the full global grade $R_U C$ and the ordered $B$-leaves.

\begin{theorem}[Causal joint terminating feedback]
Every strategy defines an exact native subprobability process
\[
 f^{\dagger,\sigma}:A\longrightarrow F_{R_U C}(B)
\]
with raw witnesses, weights, and coefficients (T). Variable node arities, zero-weight witnesses, nullary nodes, empty endpoint sets and unused global shapes are all retained as specified. No ordinary fixed-grade yanking is claimed. The policy may depend on the entire history, so this is not a state-local fixed-point equation on bare $U$: unfolding retains the history-indexed suffix strategy. Erasing that history does not in general produce one common continuation row at each $u$.
\end{theorem}
\begin{proof}
At every finite height there are countably many histories, by induction: each history has finitely many active positions and a finite product of countable witness sets is countable. Each frontier kernel is normalized. Consequently the total mass of histories terminal by height $N$, plus histories still active at height $N$, is one. This follows directly by summing each row $q_h$ and uses nonnegative countable sums only. Terminal mass is increasing with $N$ and bounded by one. Every finite terminal history appears at some height, so its total mass is their supremum and is at most one. Each terminal history has finitely many nodes, so its coefficient product exists in $\E$. The tree body has exactly the declared ordered leaves and internal state labels. No probability division or removal of raw zero histories has been used. If no terminal history exists the output raw set is empty with mass zero, but its entire grade remains $R_U C$.
\end{proof}

\paragraph{Nonhomogeneous, nonindependent worked instance.}
Take one loop state and two shapes: $b$ has arity one and returns one settled $B$-leaf, while $r$ has arity two and returns two loop occurrences. Every row chooses each shape with probability $1/2$. At each frontier choose one new fair bit and use its shape at every active position; all mixed raw tuples remain present with weight zero. A positive terminating tree has $d-1$ complete binary branching rounds followed by one all-$b$ round, probability $2^{-d}$ and $2^{d-1}$ external leaves. Summing over $d\geq1$ gives total termination mass one. The global grade is nevertheless the full $R_U C$, including every nonbalanced and mixed tree. In particular the three-leaf raw trees have weight zero here, whereas independent O1 sampling assigns their two planar forms total weight $2\cdot2^{-5}=1/16$. This is substantive joint feedback at genuinely variable output arity, not another homogeneous fragment or a mere renaming of the independent law. All finite tree coefficients remain finite products even though the expected number of executed nodes in this example is infinite.

\section{Genuine cut and assembly, with a declared boundary}
This composition combines strategies for the same declared primitive process $f$ at a retained history boundary. It is not arbitrary binary polynomial substitution of different processes or global grades, nor a newly constructed original comma bicategory. Encoding different primitives by tags would require an additional typed construction.

Fix an integer $N\geq0$ and a deterministic stopping rule on raw histories: stop by round $N$, and always stop if terminal earlier. Decisions use only the current full history. Equivalently, there is a prefix barrier $K$ meeting every path by depth $N$ exactly once, with terminal paths stopped at their terminal point. The rule and its decisions apply also to zero histories. A cut value $k\in K$ retains the \emph{entire} raw prefix tree, its settled leaves, active addressed frontier, primitive inputs, and coefficient
\[
 E_k=\prod_{v\in\mathrm{nodes}(k)}E_{x_v,w_v}.
\]
At depth zero this is the empty product. A continuation receives that entire value, not merely the leaf arity, state marginal, or list of currently active input labels. It executes all active positions jointly according to the suffix strategy; it is not a product of independent per-branch suffixes.

\begin{theorem}[Stopping-frontier cut/assembly]
For a fixed barrier, complete strategies are in bijection with a prefix strategy and one complete suffix strategy for every syntactic cut history $k$, including probability-zero histories. Their assembled terminating law satisfies
\[
 P_\sigma(t)=P_{\rm pre}(k)\,P_{{\rm suf},k}(s),\qquad
 E_t\cong E_k\times E_s,                                    \tag{C}
\]
where $t$ is uniquely reconstructed by grafting suffix forest $s$ onto $k$. Prefix mass is normalized; the suffix may be subprobabilistic. This is an actual isomorphism of the corresponding countable terminating-context objects under dependent assembly using coefficient products. At the full content-grade level assembly has the marked cut grade defined below and a directed graft/forget-cut map to $R_U C$, not generally a native process isomorphism. Repeated cuts are coherent at the retained multi-marked grade: all parenthesizations yield the same retained complete raw history, weight, and actual coefficient arrow.
\end{theorem}
\begin{proof}
Every nonterminal history lies either strictly before its unique first cut or extends one unique cut history. Thus its kernel table is assigned exactly once to prefix or suffix. Conversely those tables reconstruct every $q_h$. Condition (S) is unchanged, since suffix inputs include the full past. This proves the strategy bijection, including tables at zero histories; almost-sure quotients would not prove it.

Normalization of the prefix follows by the finite-height mass calculation with stopping made absorbing. Every terminal tree has exactly one cut prefix and one suffix forest. Splitting the finite product of its transition probabilities proves (C) literally. Splitting the finite set of node addresses gives the displayed product isomorphism in $\E$. Hence the raw bijection and its inverse preserve exact weights and carry inverse actual coefficient maps. The composite body is full tree grafting, carrying the stopping-round marker. Its global grade map is the full forget-cut map below, including unrealized marked trees and their content arrows. No claim is made that a raw-witness bijection makes this global map invertible; prefix and suffix use the retained history-indexed forest interface.

For three or more cuts, both reconstruction routes concatenate exactly the same finite list of draws and graft the same addressed tree. All coefficient composites project to the same individual node factors. Product uniqueness proves equality of the actual maps. The weight equation is associativity of a finite product; countable dependent sums agree by nonnegative Fubini. This proves cut coherence rather than postulating it as admission.
\end{proof}

\paragraph{Full marked grade, including unused cuts.}
Define $\operatorname{Cut}_U C$ to have shapes $(t,d)$ for every shape $t$ of $R_U C$ and every $0\leq d\leq\operatorname{height}(t)$. The cut after $d$ breadth-first rounds retains the prefix of nodes of depth less than $d$, all settled leaves, and an ordered forest of remaining subtrees rooted at depth $d$. At $d=0$ the prefix is one external hole and the suffix is the entire single-root tree; at $d=\operatorname{height}(t)$ the suffix is the empty forest. The content coefficient is the tensor of all node factors, or equivalently the tensor of prefix factors and forest factors with their canonical addressed permutation. All such cuts are included, whether or not the actual stopping rule realizes them. The full arity-preserving map
\[
 \gamma:\operatorname{Cut}_U C\longrightarrow R_U C,
 \qquad(t,d)\longmapsto t
\]
has the actual canonical coefficient arrow. The permissible partial trees are precisely level-complete prefixes: every node strictly above level $d$ is present, each already settled external leaf is retained, and the holes are exactly the still-active loop edges reaching level $d$. The ordered suffix consists of one nonempty rooted tree at each such hole, with the matching loop label; the special $d=0$ case has the single external root hole. A terminal prefix permits no suffix and has $d$ equal to its actual height. Thus arbitrary partial trees with holes at unequal depths, arbitrary forests, or trailing idle cut rounds are not admitted. Grafting these typed partial trees and forests is bijective onto this \emph{marked} grade, not onto the unmarked grade. For every nonempty tree the two distinct markers $0$ and $\operatorname{height}(t)$ already show that $\gamma$ is generally noninjective. The raw cut body chooses $d$ from the stopping history and maps to this full grade. Thus the cut theorem gives an invertible context map paired with a generally noninvertible global content map. If an isomorphism of full processes is required, both sides must retain this same marker. Two successive cuts use all $(t,d_1,d_2)$ with $0\leq d_1\leq d_2\leq\operatorname{height}(t)$; reassociation is an isomorphism at that full multi-marked interface. No unused cut is discarded. The suffix object is explicitly a forest over the retained frontier, including the empty forest, rather than another single-root $R_U C$ process.

The cut theorem is not an assertion that arbitrary independently supplied branch laws have a causal common coupling. Its suffix is a joint strategy for the whole retained frontier. In particular a cut may not separate simultaneous sibling draws and then silently treat one sibling outcome as prior history while retaining the same strategy interface.

\section{Actual native comparisons and the old O1 injection}
Suppose primitive processes $f,f'$ have a native comparison: maps $a:A\to A'$, $u:U\to U'$, $b:B\to B'$, one full arity-preserving collection map $\alpha:C\to C'$, and rowwise raw maps $r_x:W_x\to W'_{x'}$ preserving probabilities and actual context coefficients. The body equations preserve each ordered address, its shape image and whether it terminates or continues. Consequently there is a map $H$ on every raw history and a coordinatewise map $r_h$ on its frontier tuples. Declare a strategy comparison by the local equations
\[
 (r_h)_*q_h=q'_{Hh}\quad\hbox{for every source history }h.     \tag{M}
\]
This is an explicit conditional-kernel compatibility (a lumpability condition), not a whole-execution comparison assumed in advance. In particular different histories identified by $H$ must give the same target conditional law. General shared strategies need not satisfy it.

\begin{theorem}[Native transport]
Under (M), the induced raw terminal-tree map has exact fibre pushforward. Its nodewise primitive coefficient arrows induce actual product coefficient arrows, and the global grade map is the full $R_u\alpha$, with all unused shapes and their tensor coefficients. These comparisons preserve identities and composition. They commute with cut assembly whenever the two stopping rules commute with $H$.
\end{theorem}
\begin{proof}
The statement is rowwise for each specified external input and its image under the displayed external-input map; it does not sum the process rows over all external roots having that image. Induct on height. For a target prefix, sum over source prefixes above it and over source next tuples above its next tuple. Condition (M) makes the inner sum the same target transition probability at each source prefix. Multiplication by the inductive prefix mass gives the target next-prefix mass. The base root is sent to its specified root, of mass one. Arity and termination preservation ensure every fibre lies in the same height and active/terminal stratum. Summing these exact equalities over terminal heights proves exact terminal pushforward. This proof also covers zero target masses without choosing conditional distributions by division.

At a target terminal tree, project the source coefficient product to each corresponding source node factor and apply that primitive actual coefficient arrow. The ordered node bijection gives an arrow into the target product. Full content functoriality gives $R_u\alpha$, using the actual arrows at every node of every global tree, not only realized trees. Product uniqueness and tensor coherence prove identity and composition equations. Compatible cuts simply partition the same node factors and finite list of transition kernels, so the comparison commutes with (C).
\end{proof}

The distinguished independent strategy is
\[
 q^{\rm ind}_h(\mathbf w)=\prod_{i\in I(h)}p_{x_i(h)}(w_i).
\]
Every original rowwise native comparison satisfies (M) between these strategies: the fibre sum of a finite product factors into the product of the primitive fibre sums. Thus \emph{all} original arity-preserving grade and actual coefficient comparisons transport on the independent locus, including arbitrary same-arity shape merges. This does not imply that every shared strategy transports along every such merge.

The original independent recursive construction has the same raw trees and weights at this distinguished strategy, but context coefficient $\bigotimes_v E_{x_v,w_v}$. There is a canonical actual native comparison
\[
 f^{\dagger,\mathrm{old}}\longrightarrow f^{\dagger,\sigma_{\rm ind}},
 \qquad \bigotimes_v E_v\longrightarrow\prod_v E_v,            \tag{O1}
\]
using each tensor projection obtained by augmenting all other factors. It is identity on raw trees and on the entire global grade $R_U C$. These maps are natural in all original native arrows and coherent with cuts: their projection to any node is that node's map with all other tensor factors augmented. Product uniqueness proves the equality, so no faithful-carrier premise is used. Finite vector-space coefficients make (O1) generally noninvertible. Here cut naturality is understood at the marked grade on both sides before the directed forgetting map $\gamma$; it does not assert an unmarked cut isomorphism. The collection of all strategies can be retained as an index family of these subprobability objects; (O1) selects the independent component, not a canonical coupling of every strategy.

\section{Decisive boundaries and countermodels}
\begin{proposition}[Scope-splitting obstruction]
No interface can simultaneously allow all within-scope couplings satisfying (S), admit every sequential cut through a scope while revealing the first outcome, and require the unchanged primitive marginal for the second draw conditional on the revealed history.
\end{proposition}
\begin{proof}
A scope with two fair-bit rows admits $q(0,0)=q(1,1)=1/2$. After a proposed cut revealing the first bit $X$, the second bit $Y$ is deterministically $X$. Its conditional fair-bit marginal fails at both positive histories. This is a failure of the required law, not of coefficients or a zero-history convention. It embeds with terminal coefficients and fixed binary shape in the original O1 input, hence also in the variable-arity setting.
\end{proof}

A second obstruction distinguishes branchwise history from global history. Let two branches each have two fair bits $(X_i,Y_i)$ with independent within-branch law. Jointly choose independent fair $X_1,X_2$ and set $Y_1=X_2$, $Y_2=X_1$. Both branch trajectories have their prescribed independent laws, yet after the first round the second-round primitive bits are deterministic given the global history. Consequently arbitrary couplings of complete branch trajectories need not define causal frontier strategies. The retained cut theorem does not authorize this coupling.

The earlier variable-arity routing obstruction likewise fails admission here: an always-active continuation slot was set equal to a previously revealed routing bit. Its conditional marginal given the routing history is not fair. By contrast, a new fair bit shared diagonally by every active continuation in that round is admitted. Its law remains fair conditional on the routing history while preserving nontrivial sibling correlation. Thus the construction repairs the routing failure without retreating to independent occurrence tensors or homogeneous grades.

\paragraph{What remains open.}
The theorem establishes a scope-atomic causal strategy doctrine with full variable-arity recursive grades and exact cut assembly. It does not establish an unmarked polynomial-substitution bicategory for arbitrary shared laws, invariance under every scheduler or scope refinement, arbitrary persistent mutable interference, or an infinite-stream semantics. More permissive cuts require a residual scope object recording a conditional joint law and changing its advertised primitive kernels; the fair-bit obstruction proves that such extra state cannot be erased. The required dependent residual-law compiler, its native coefficient semantics, and its compatibility with unrestricted original grade maps are separate obligations, not conclusions of the cut theorem.
\end{document}
